% Project-wide additions loaded after the LiX-based textbook class.

\usepackage[most]{tcolorbox}
\usepackage{tikz}
\usepackage{eso-pic}

% Short corner marks at the four corners of the text block, matching the
% textbook.cls geometry (top=30mm, bottom=30mm, left=24mm, right=24mm). A
% restrained, print-shop-style cue for where the margin begins, activated
% from Part I onward in main.tex so the bespoke cover, imprint, table of
% contents, and roadmap pages stay exactly as designed.
\newcommand{\pagecornermarks}{%
    \begin{tikzpicture}[overlay, line width=0.5pt, color=black!55]
        \draw ([xshift=24mm,yshift=-30mm]current page.north west) -- ++(8mm,0);
        \draw ([xshift=24mm,yshift=-30mm]current page.north west) -- ++(0,-8mm);
        \draw ([xshift=-24mm,yshift=-30mm]current page.north east) -- ++(-8mm,0);
        \draw ([xshift=-24mm,yshift=-30mm]current page.north east) -- ++(0,-8mm);
        \draw ([xshift=24mm,yshift=30mm]current page.south west) -- ++(8mm,0);
        \draw ([xshift=24mm,yshift=30mm]current page.south west) -- ++(0,8mm);
        \draw ([xshift=-24mm,yshift=30mm]current page.south east) -- ++(-8mm,0);
        \draw ([xshift=-24mm,yshift=30mm]current page.south east) -- ++(0,8mm);
    \end{tikzpicture}%
}

% Two accents serve the entire book: orange for chapter-level outcomes and
% summaries, and a restrained blue for navigation, code, and study prompts.
\definecolor{BookBlue}{RGB}{30,78,112}
\definecolor{BookOrange}{RGB}{220,112,0}

% A compact callout used throughout the self-study chapters.
\newtcolorbox{NoteBox}{
    breakable,
    colback=black!1,
    colframe=BookBlue,
    boxrule=0.6pt,
    arc=1.5mm,
    left=2mm,
    right=2mm,
    top=1.5mm,
    bottom=1.5mm
}

% Reusable learning components. Each has one job so chapters can add depth
% without becoming a wall of undifferentiated prose.
\newtcolorbox{SyntaxBox}[1]{
    breakable,
    title={Syntax focus: #1},
    colback=black!1,
    colframe=BookBlue,
    fonttitle=\bfseries,
    boxrule=0.7pt,
    arc=1.5mm
}

\newtcolorbox{ExplainBox}[1]{
    breakable,
    title={Why it works: #1},
    colback=black!1,
    colframe=BookBlue,
    fonttitle=\bfseries,
    boxrule=0.7pt,
    arc=1.5mm
}

\newtcolorbox{TryBox}[1]{
    breakable,
    title={Hands-on checkpoint: #1},
    colback=black!1,
    colframe=BookBlue,
    fonttitle=\bfseries,
    boxrule=0.7pt,
    arc=1.5mm
}

% Centralize image inclusion. The third argument stores a concise description
% that is also printed as the figure caption by \bookfigure. Do not pass the
% experimental graphicx `alt` key in the standard build: it activates tagging
% code that is not yet reliable with this book's listings and boxed content.
\newcommand{\accessibleimage}[3]{%
    \includegraphics[#1]{#2}%
}

% Standard figure treatment for both existing illustrations and generated
% program outputs. Every instructional graphic receives a descriptive caption.
\newcommand{\bookfigure}[4][0.92\textwidth]{%
    \begin{figure}[htbp]
        \centering
        \accessibleimage{width=#1}{#2}{#3}
        \caption{#3}
        \label{#4}
    \end{figure}
}

% Listings uses named languages. Define the two lightweight formats used for
% terminal commands and plain output so every chapter compiles consistently.
\lstdefinelanguage{text}{}
\lstdefinelanguage{bash}{
    sensitive=true,
    morecomment=[l]{\#},
    morestring=[b]"
}

% LiX implements \c with \lstinline, which is verbatim-like and therefore
% fails inside bold labels, list items, and other command arguments. This
% protected, detokenized form preserves code punctuation in those contexts.
\RenewDocumentCommand{\c}{m}{\texttt{\detokenize{#1}}}

% Keep captions with the code blocks they describe and place them underneath,
% matching the original \code command's layout. The restrained syntax palette
% uses one strong accent for language keywords, one supporting color for
% strings, and neutral gray for comments; numbers and ordinary names stay
% black for print readability.
\colorlet{PythonKeyword}{BookBlue}
\definecolor{PythonString}{RGB}{38,105,82}
\definecolor{PythonComment}{RGB}{100,100,100}

\lstdefinestyle{block}{
    basicstyle=\ttfamily\footnotesize,
    keywordstyle=\color{PythonKeyword}\bfseries,
    stringstyle=\color{PythonString},
    commentstyle=\color{PythonComment}\itshape,
    breaklines=true,
    numbers=left,
    numberstyle=\scriptsize\color{black!55},
    numbersep=7pt,
    aboveskip=0.5em,
    belowskip=0em,
    columns=fullflexible,
    keepspaces=true,
    upquote=true,
    showstringspaces=false,
    xleftmargin=14pt,
    postbreak=\mbox{\hspace{-2.5em}\textcolor{gray}{$\hookrightarrow$}\space\space}
}

\lstset{captionpos=b}

% LiX hides link styling. Visible but restrained links make the table of
% contents and cross-references discoverable on screen while remaining clean
% in print.
\hypersetup{
    colorlinks=true,
    linkcolor=BookBlue,
    urlcolor=BookBlue,
    citecolor=BookBlue,
    linktoc=all,
    bookmarksopen=true,
    bookmarksnumbered=true,
    pdflang={en-US},
    pdfstartview=FitH
}

% Preserve Unicode mappings in engines that expose this PDF primitive.
\ifdefined\pdfgentounicode
    \pdfgentounicode=1
\fi

\newcommand{\booklink}[2]{\hyperref[#1]{#2}}
